%% Coin — AI Study Planner for Grad Students
%% IEEE conference format (IEEEtran).
%% Compile on Overleaf (IEEEtran is built in) or locally:
%%   pdflatex report.tex   (run twice for references)
\documentclass[conference]{IEEEtran}
\IEEEoverridecommandlockouts

\usepackage{cite}
\usepackage{amsmath,amssymb}
\usepackage{graphicx}
\usepackage{booktabs}
\usepackage{array}
\usepackage{url}
\usepackage[hidelinks]{hyperref}

\begin{document}

\title{Coin: A Private, Note-Grounded AI Study Planner\\ for Graduate Students}

\author{\IEEEauthorblockN{Seohee Han}
\IEEEauthorblockA{\textit{Generative AI and Blockchain} \\
\textit{Gwangju Institute of Science and Technology (GIST)}\\
Gwangju, South Korea \\
hanseohee@gm.gist.ac.kr}}

\maketitle

\begin{abstract}
Graduate research removes the explicit yardstick of undergraduate exams: there
is no weekly score, advisor feedback is coarse, and general-purpose chatbots are
stateless strangers that return plausible-sounding generalities because they have
no memory of the student. This paper presents \textbf{Coin}, a private Personal AI
Butler that runs entirely on a personal laptop. Coin wires \emph{Claude} into a
graduate student's real tools---Obsidian notes, a schedule store, a music log, and
a blog---through five self-authored Model Context Protocol (MCP) servers, and is
reached conversationally through Telegram. Its defining capability is a weekly
research review produced by Retrieval-Augmented Generation (RAG): it retrieves the
week's journal, the student's research-goal file, and an activity summary, then
generates feedback \emph{grounded in those records} rather than inventing them. We
analyze the system through the agent-intelligence frame $I = M \times
\mathrm{HBM} \times R$ and show non-trivial mass in all three factors, document the
cost and rate-limit engineering that makes daily use affordable, and present an
11-day usage log as practicality evidence.
\end{abstract}

\begin{IEEEkeywords}
AI agent, Model Context Protocol, retrieval-augmented generation, personal
assistant, local-first, cost-efficient LLM.
\end{IEEEkeywords}

\section{Introduction}
The transition from undergraduate to graduate study removes a feedback signal that
students rarely notice until it is gone. Exams provided a recurring, quantitative
answer to \emph{``am I doing okay?''}. Research provides none. The weekly lab
meeting is the only structured checkpoint, and it is necessarily coarse---an
advisor cannot give fine-grained, day-level feedback, and a student cannot report
everything.

Three existing tools fail to fill this gap. \emph{Paper planners} are tedious to
maintain, are designed for exam-driven study, and give no feedback. \emph{General
chatbots} (ChatGPT, Gemini) have no persistent memory of the student's goals; each
session starts from zero, so advice is generic and nothing accumulates.
\emph{Note apps} (Obsidian, Notion) store records faithfully but do not reason over
them.

We argue that the missing artifact is a \emph{personalized planner that (a)
remembers the student and (b) gives weekly feedback grounded in the student's own
records.} This paper describes such a system, Coin, built as a private AI butler
on the OpenClaw + Claude + MCP stack.

\section{System Architecture}
Coin runs on an LG gram 16 (Intel i7-1195G7, 16~GB RAM) under Ubuntu 24.04 LTS.
The OpenClaw gateway runs continuously as a \texttt{systemd} service, so the
assistant is always reachable. The user interface is a Telegram bot; the reasoning
engine is the Claude API; the connection to the user's data is MCP.

Conceptually, the OpenClaw gateway (``Coin'') sits at the center: Telegram provides
the chat interface, the Claude API (Sonnet/Haiku) provides reasoning, and five
local MCP servers connect outward to the Obsidian Vault, a Notion mirror, a Jekyll
blog, a Spotify listening log, and read-only web access.

Five local MCP servers expose the user's tools to Claude
(Table~\ref{tab:mcp}). None stores credentials in source; each reads secrets from
\texttt{\textasciitilde/.openclaw/*.json} at runtime.

\begin{table}[t]
\caption{MCP servers and their tools}
\label{tab:mcp}
\centering
\footnotesize
\begin{tabular}{@{}llp{4.6cm}@{}}
\toprule
\textbf{Server} & \textbf{Tools} & \textbf{Function} \\
\midrule
\texttt{brain}    & 13 & Read / search / write the Obsidian Vault \\
\texttt{schedule} & 17 & Calendar, todos, commitments, \texttt{weekly\_summary} \\
\texttt{spotify}  & 15 & Listening log, ``obsessed-on'', playback, heatmap \\
\texttt{blog}     & 6  & Publish posts/news to a Jekyll GitHub Pages blog \\
\texttt{usage}    & 5  & Token/cost/time tracking from session files \\
\texttt{fetch}    & -- & Read-only web access (\texttt{mcp-server-fetch}) \\
\bottomrule
\end{tabular}
\end{table}

\section{Core Mechanism: Note-Grounded Weekly Review}
A weekly review is the one task a generic chatbot cannot perform, because the
required context---\emph{this student's} week---does not exist in its weights. Coin
treats the review as a Retrieval-Augmented Generation problem.

\textbf{Retrieval.} Before any text is generated, Coin pulls three sources:
(1)~\texttt{brain.recent\_logs(7)}---the week's seven journal entries;
(2)~\texttt{brain.read\_note(\textnormal{``RESEARCH\_GOALS''})}---the student's goal
file; and (3)~\texttt{schedule.weekly\_summary}---the week's activity mix.

\textbf{Generation.} Claude composes the review \emph{grounded in} that retrieved
context. In the demo, the review correctly states that the student ``finished the
DAE paper,'' prepared ``the lab-meeting PPT,'' and ``slept at 5~AM''---each fact
retrieved from a journal entry, not hallucinated. A plain chatbot, lacking these
records, instead produces plausible-sounding generalities. This contrast is the
project's central empirical claim.

\section{Analysis Frame: $I = M \times \mathrm{HBM} \times R$}
Course material framed LLM-substrate intelligence as $I = \mathrm{Compute} \times
\mathrm{Memory} \times \mathrm{Retrieval}$. We re-map this at the \emph{agent
layer}:
\begin{equation}
I_{\mathrm{agent}} = M \times \mathrm{HBM} \times R
\end{equation}
where $M$ is the curated \texttt{.md} state (the student's accumulating Obsidian
notes), $\mathrm{HBM}$ is the chosen model and loaded context, and $R$ is the
MCP/RAG/tool layer. The product is \emph{multiplicative}: any factor near zero
collapses agent intelligence, so a useful butler must show non-trivial mass in all
three.

\textbf{$M$ (curated state)---non-trivial.} Coin's memory is not a chat buffer but
a structured, growing Vault: daily journal entries, a \texttt{RESEARCH\_GOALS.md}
goal file, Brain notes (books, papers, insights), and saved weekly reviews. Over
the evaluation window the journal alone accumulated 11 days of tagged, structured
entries. This is durable, queryable state, not ephemeral context. \emph{Failure
mode avoided:} with low $M$ the butler would be stateless and forget the user.

\textbf{$\mathrm{HBM}$ (model / context)---non-trivial.} Coin deliberately routes
work across a two-tier model stack: Haiku~4.5 for everyday calls and Sonnet~4.6
reserved for the weekly synthesis where reasoning quality matters. Context is
actively shaped---\texttt{lightContext} trims persona and background to the
minimum, and \texttt{toolsAllow} exposes only the ${\sim}6$ tools a weekly review
needs instead of all ${\sim}50$. \emph{Failure mode avoided:} with low
$\mathrm{HBM}$ the butler could not load consolidated memory within usable context.

\textbf{$R$ (retrieval / tools)---non-trivial.} Five MCP servers expose ${\sim}50$
tools, and the weekly review exercises a genuine RAG pipeline that retrieves from
the Vault and schedule store before generating. Retrieval is observable in the
usage log as concrete tool calls with outcomes and latencies. \emph{Failure mode
avoided:} with low $R$ the butler would work from stale data and hallucinate the
present.

Because all three factors carry real mass, the product---useful agent
intelligence---does not collapse.

\section{Cost and Rate-Limit Engineering}
\textbf{The rate-limit wall.} The Anthropic API Tier-1 limit is 30{,}000 input
tokens per minute. A naive agent re-sends, on every step, both the entire
conversation so far and the tool manuals for all ${\sim}50$ tools---so even ``read
one note'' can exceed 30k and the API rejects the request, like re-reading 50
manuals before every errand.

\textbf{Fix: shrink what gets sent.} Three measures, in order of impact:
(1)~\texttt{toolsAllow}---expose only the ${\sim}6$ tools a weekly review needs and
drop the other ${\sim}44$ manuals (largest saving);
(2)~\texttt{lightContext}---trim persona and background context to the minimum;
(3)~model choice---Sonnet only for the weekly synthesis, the cheap model
otherwise. Requests then fit under 30k, the rate-limit errors disappear, and
per-call cost drops. Building the agent \emph{light} yields stability and savings
at once.

\textbf{Cost economy.} Because the system bills per token rather than by
subscription, the design question is not ``what \emph{can} it do?'' but ``what is
\emph{worth paying for}?'' Capabilities that free local software already provides
were deliberately subtracted (Table~\ref{tab:subtraction}). What remains is an
on-demand tool plus a single weekly cron job---the one task uniquely Coin's. The
\texttt{usage} MCP measures tokens, cost, and active time, so spend is monitored,
not guessed; the estimated steady-state cost of the final design is on the order of
a few US dollars per month for personal daily use.

\begin{table}[t]
\caption{Subtraction: keep only what is uniquely worth paying for}
\label{tab:subtraction}
\centering
\footnotesize
\begin{tabular}{@{}llp{3.0cm}@{}}
\toprule
\textbf{Capability} & \textbf{Decision} & \textbf{Reason} \\
\midrule
Sending email          & removed   & I can send it myself \\
Spotify / Notion sync  & removed   & Linux cron does it, free \\
Daily song alerts      & removed   & just open Obsidian / Notion \\
Blog edits             & removed   & Claude Code does it better \\
\textbf{Weekly research feedback} & \textbf{kept} & \textbf{only Coin can do this} \\
\bottomrule
\end{tabular}
\end{table}

\section{Privacy and Security}
\textbf{Local-first data flow.} All personal data is created and stored on the
laptop (SQLite + Obsidian markdown) and never leaves it except as the minimal
context a single Claude request requires. Sources on the laptop---Obsidian
\texttt{.md}, \texttt{schedule.db}, \texttt{spotify.db}---are read by a local MCP
server, which forwards only a minimal per-request slice to the cloud Claude API;
the grounded answer is then written back to the Vault. The bulk corpus stays local.

\textbf{Threat model.} Table~\ref{tab:threat} summarizes assets, threats, and
mitigations. There were zero security incidents during the project period: no
credential was committed and no personal data left the device beyond per-request
context.

\begin{table*}[t]
\caption{Threat model}
\label{tab:threat}
\centering
\footnotesize
\begin{tabular}{@{}lll@{}}
\toprule
\textbf{Asset} & \textbf{Threat} & \textbf{Mitigation} \\
\midrule
API / bot keys & Leak via source repo & Secrets only in \texttt{\textasciitilde/.openclaw/*.json}; \texttt{.gitignore} excludes them; repo scanned for hard-coded keys (none found) \\
Personal records (journal, spend, listening) & Exposure in public repo & All \texttt{*.db} and Vault content excluded from the repo \\
Unauthorized bot use & Stranger messages the bot & Telegram restricted to a single owner ID; group messages require explicit mention \\
Cloud over-exposure & Whole Vault sent to API & \texttt{lightContext} + \texttt{toolsAllow} send only the slice a request needs \\
\bottomrule
\end{tabular}
\end{table*}

\section{Differentiation vs. Big-Tech Assistants}
Coin demonstrates at least three OpenClaw advantages over hosted assistants:
(1)~\emph{local \& private}---runs on my laptop; records never centralize in a
vendor cloud; (2)~\emph{note-grounded (RAG)}---feedback cites my real records
instead of inventing generalities; (3)~\emph{accumulating memory}---the Vault
grows, so each week's review is better informed than the last;
(4)~\emph{self-extensible via MCP}---five servers I wrote connect Claude to
Obsidian, Notion, Spotify, my blog, and the web; (5)~\emph{cost-bounded}---a
two-tier model stack and trimmed context keep daily use under the rate limit and
affordable.

\section{Evaluation}
\textbf{Practicality.} An 11-day usage log records MCP-mediated calls in the form
\texttt{HH:MM CHANNEL server task (OUTCOME, latency)} across two channels (CLI and
Telegram); see Table~\ref{tab:usage}. It evidences daily, real use rather than a
one-off demo, and shows the retrieval calls underlying the weekly review. (The
course requires a 7-day log; 11 days are provided.)

\begin{table}[t]
\caption{11-day usage log summary (2026-05-20 to 2026-06-07)}
\label{tab:usage}
\centering
\footnotesize
\begin{tabular}{@{}ll@{}}
\toprule
\textbf{Metric} & \textbf{Value} \\
\midrule
Days covered     & 11 \\
Total calls      & 206 (203 OK, 3 FAILED; 98.5\% success) \\
Channels         & CLI 201; Telegram 5 \\
Busiest servers  & \texttt{brain} 97; \texttt{schedule} 81; \texttt{spotify} 27 \\
\bottomrule
\end{tabular}
\end{table}

\textbf{Qualitative.} In the recorded demo the weekly review reproduces
week-specific facts (DAE-paper completion, lab-meeting PPT, a 5~AM sleep night, and
a 23{,}900-KRW stopwatch purchase flagged as the week's largest single expense),
none of which a memoryless chatbot could supply.

\section{Future Work}
First, \emph{counsel first, then feedback}: before writing a review, ask
${\sim}5$ short questions over Telegram (``Why didn't the paper get read?'', ``What
felt most stuck this week?'') and fold the answers into the feedback, turning a
one-way review into coaching. Second, \emph{smarter as records stack}: as weeks
accumulate into months and quarters, surface longitudinal patterns (``this month
vs. last quarter'', ``what keeps getting pushed back?'') and let the review refine
\texttt{RESEARCH\_GOALS.md} over time.

\section{Conclusion}
Coin shows that a useful private AI butler for a graduate student is not a matter
of raw model capability but of getting three multiplicative factors all
non-trivial at once: a curated, accumulating note state ($M$); deliberate model and
context selection ($\mathrm{HBM}$); and a real MCP/RAG retrieval layer ($R$).
Grounding generation in the student's own records turns a generic chatbot into a
planner that remembers, and disciplined cost and context engineering makes it cheap
enough to use every day.

\begin{thebibliography}{1}
\bibitem{openclaw} OpenClaw documentation. [Online]. Available: \url{https://openclaw.ai}
\bibitem{anthropic} Anthropic, ``Claude API and Model Context Protocol (MCP).'' [Online]. Available: \url{https://docs.anthropic.com}
\bibitem{lee} H.-N. Lee, ``Building a Useful Private AI Butler,'' GIST INFONET Lab, course notes (Week 10: \emph{Intelligence $=$ Compute $\times$ Memory $\times$ Retrieval}), 2026.
\end{thebibliography}

\end{document}
